\pdfvariable suppressoptionalinfo 611
\providecommand{\CRevisionRound}{2}
\documentclass[10pt]{article}
\usepackage[a4paper,margin=17mm]{geometry}
\usepackage{amsmath,amssymb,amsthm,booktabs,microtype}
\usepackage[T1]{fontenc}
\usepackage{lmodern}
\usepackage[hidelinks,hypertexnames=false]{hyperref}
\newtheorem{theorem}{Theorem}
\newtheorem{proposition}{Proposition}
\title{A Coordinate-Safe Phase Portrait for the Bose--Josephson Dimer}
\author{HCS Research Program}
\date{30 August 2026\quad (revision \CRevisionRound)}
\begin{document}
\maketitle
\begin{abstract}
We study the normalized two-mode Bose--Josephson dimer
\(H(z,\phi)=\Lambda z^2/2-\sqrt{1-z^2}\cos\phi\), \(\Lambda\ge0\).
Passing to the Bloch sphere removes the apparent coordinate singularity at
\(z=\pm1\).  We classify \((0,0)\), \((0,\pi)\), and the
\(\Lambda>1\) symmetry-broken equilibria, including the \(\Lambda=1\)
pitchfork.  Elimination of the phase gives an exact quartic in \(z\); its
roots yield complete-elliptic-\(K\) periods on crossing and self-trapped
components.  The \(H=1\) sech homoclinic, its \(\Lambda=2\) pole limit, and
the reverse component criterion are stated without assigning a finite period
to a separatrix.  This is a source-local phase-portrait theorem: no arithmetic
owner or target determinant is claimed.
\end{abstract}

\section{Hamiltonian and Bloch regularization}
Set
\[
 H(z,\phi)=\frac{\Lambda z^2}{2}-\sqrt{1-z^2}\cos\phi,
 \qquad \Lambda\ge0,
\]
with canonical equations
\[
 \dot z=-\sqrt{1-z^2}\sin\phi,\qquad
 \dot\phi=\Lambda z+\frac{z\cos\phi}{\sqrt{1-z^2}}. \tag{1}
\]
The Bloch variables
\(x=\sqrt{1-z^2}\cos\phi\), \(y=\sqrt{1-z^2}\sin\phi\) satisfy
\[
 \dot x=-\Lambda zy,\qquad \dot y=z(1+\Lambda x),\qquad \dot z=-y. \tag{2}
\]
Thus the north and south poles have \(\dot y=+1\) and \(-1\), respectively;
they are not equilibria and no value of the singular angle \(\phi\) is
assigned there.

\begin{theorem}[Fixed points and pitchfork]
The points \((0,0)\) and \((0,\pi)\) exist for every \(\Lambda\).  The first
has energy \(-1\), linearization
\(\left[\begin{smallmatrix}0&-1\\\Lambda+1&0\end{smallmatrix}\right]\),
and frequency \(\sqrt{\Lambda+1}\).  The second has energy \(+1\) and
linearization \(\left[\begin{smallmatrix}0&1\\\Lambda-1&0\end{smallmatrix}\right]\):
it is elliptic for \(\Lambda<1\), parabolic at \(\Lambda=1\), and hyperbolic
for \(\Lambda>1\).  For \(\Lambda>1\), two additional elliptic points are
\[
 z=\pm\sqrt{1-\Lambda^{-2}},\quad \phi=\pi,\quad
 H_{\max}=\frac{\Lambda+\Lambda^{-1}}2. \tag{3}
\]
In the canonical \((z,\phi-\pi)\) coordinates their matrix is
\(\left[\begin{smallmatrix}0&1/\Lambda\\-\Lambda(\Lambda^2-1)&0\end{smallmatrix}\right]\),
so their small-amplitude period is
\(2\pi/\sqrt{\Lambda^2-1}\).  The zero-phase crossing limit is
\(2\pi/\sqrt{\Lambda+1}\), and (3) coalesces with \((0,\pi)\) at the
pitchfork \(\Lambda=1\).
\end{theorem}
\begin{proof}
Differentiate (1) and solve \(\dot z=\dot\phi=0\) away from the poles.
The displayed matrices follow by first-order expansion.  Equation (2) is
obtained by the chain rule and extends the vector field over both poles.
\end{proof}

\section{Energy quartic and regular periods}
On \(H=h\),
\[
 \dot z^2=(1-z^2)-\left(\frac{\Lambda z^2}{2}-h\right)^2
 =-\frac{\Lambda^2}{4}z^4+(\Lambda h-1)z^2+1-h^2. \tag{4}
\]
For \(\Lambda>0\), let
\[
 y_\pm=\frac{2(\Lambda h-1\pm
 \sqrt{\Lambda^2-2\Lambda h+1})}{\Lambda^2}. \tag{5}
\]
Then \(\dot z^2=(\Lambda^2/4)(y_+-z^2)(z^2-y_-)\).  If
\(-1<h<1\), \(y_-<0<y_+\), the level is one connected crossing component and
\[
 T_{\rm cross}=\frac{8}{\Lambda\sqrt{y_+-y_-}}
 K\!\left(\sqrt{\frac{y_+}{y_+-y_-}}\right). \tag{6}
\]
If \(\Lambda>1\) and \(1<h<H_{\max}\), there are two sign-preserving
components and each has
\[
 T_{\rm self}=\frac{4}{\Lambda\sqrt{y_+}}
 K\!\left(\sqrt{1-\frac{y_-}{y_+}}\right). \tag{7}
\]
Here \(K\) uses the modulus convention; the code passes its square to the
numerical library.  The limits near the two elliptic centers are exactly the
frequencies stated above.

\section{Separatrix and component criterion}
For \(\Lambda>1,h=1\), (4) has the homoclinic solution
\[
 z(t)=\pm\frac{2\sqrt{\Lambda-1}}{\Lambda}
 \operatorname{sech}(\sqrt{\Lambda-1}\,t). \tag{8}
\]
The full critical level is connected through the saddle at \(z=0\), with two
one-sided homoclinic branches.  At its turning point the phase is \(\pi\) for
\(1<\Lambda<2\), the point is a Bloch pole at \(\Lambda=2\), and the phase is
0 for \(\Lambda>2\).  Formula (8) has infinite period; it is not a regular
periodic orbit.

For regular \(\Lambda>1\) levels, \(1<h<H_{\max}\) is exactly the
self-trapped regime: the two connected components preserve the sign of \(z\),
and changing the initial sign selects the reflected component.  For
\(-1<h<1\), the single component crosses \(z=0\), so both signs occur.  At
\(h=1\), sign change is only asymptotic.  At \(\Lambda=1,h=1\), (4) reduces
to \(-z^4/4\), so only the isolated degenerate point \(z=0\) remains.  At
\(\Lambda=0\), \(H=-x\) gives rigid Bloch rotation of period \(2\pi\) on
regular circles.

\ifnum\CRevisionRound>0
\section{Certified finite receipt}
The accompanying receipt contains 14 fixed-point rows, 8 pole rows, 13 level
rows, and 5 component criteria.  A producer-independent checker makes 995
assertions.  SymPy verifies the Hamilton/Bloch identities, quartic roots,
pitchfork and sech profile, while three independent transformed quadratures
reproduce (6)--(7).  Byte replay passes and a hostile suite rejects 28/28
repaired-hash mutations.
\fi

\ifnum\CRevisionRound>1
\section{Route-A boundary}
Regular levels form a continuum, so the result is A1\_WEAK rather than a
discrete primitive-orbit atlas.  There is no intrinsic rational-prime carrier,
target weighted zeta, Fredholm determinant, or target-zero match.  The locked
tuple is
\[
(\mathtt{A0\_FAIL},\mathtt{A1\_WEAK},\mathtt{A2\_FAIL},
 \mathtt{A3\_FAIL},\mathtt{A4\_NATURAL\_QUANTIZATION}),
\]
with \texttt{ROUTE\_A\_REJECTED}, Route B disabled, and scope literal
\texttt{NO\_BAD\_EULER\_OR\_ROOT\_NUMBER}.  Source-local phase portraits
must not be reinterpreted as target arithmetic structures.
\fi

\begin{thebibliography}{9}
\bibitem{smerzi} A. Smerzi, S. Fantoni, S. Giovanazzi and S. Shenoy,
Quantum coherent atomic tunneling between two trapped Bose--Einstein
condensates, \emph{Phys. Rev. Lett.} 79 (1997), 4950--4953,
DOI: 10.1103/PhysRevLett.79.4950, \url{https://doi.org/10.1103/PhysRevLett.79.4950}.
\bibitem{raghavan} S. Raghavan, A. Smerzi, S. Fantoni and S. Shenoy,
Coherent oscillations between two weakly coupled Bose--Einstein condensates,
\emph{Phys. Rev. A} 59 (1999), 620--633, DOI: 10.1103/PhysRevA.59.620,
\url{https://doi.org/10.1103/PhysRevA.59.620}.
\end{thebibliography}

\section*{Declarations}
\textbf{Scope.} \texttt{NO\_BAD\_EULER\_OR\_ROOT\_NUMBER}; no target
arithmetic, target-zero, determinant-matching, or Hilbert--Pólya claim.
\textbf{Data and code.} All rows and audits accompany HCS-C243. This is not
external peer review. \textbf{AI-use disclosure.} Generative tools assisted
drafting and code generation; the artifact chain checks displayed formulas and
metadata.
\end{document}
